\section{Introduction}
This text concerns the use of computational problems over lattices as a way to construct cryptography. The introduction of lattices into cryptography has been a fruitful endeavour. In the terminology of Impagliazzo \cite{imp95}, lattices started in Minicrypt with Ajtai's introduction of the Short Integer Solutions (SIS) problem, allowing the construction of one-way functions. With time, however, we can definitely say lattices have entered Cryptomania, where public-key encryption and more advanced primitives are possible. As things stand today, cryptography from lattices serves as the foundation for standarised cryptography in the presence of quantum computers, which motivates a solid understanding of the mathematics of, computational problems on, and cryptographic primitives from, lattices. 

The attractiveness of lattice-based cryptography is many-fold. We present three reasons from \cite{peikert_dec}.
\begin{itemize}
    \item As mentioned above, lattices have proven a useful tool for realising many of the goals of cryptography. Apart from public-key encryption and key encapsulation mechanisms, lattices have been used to construct homormorphic encryption schemes, attribute-based encryption and general-purpose code obfuscation. Homormorphic encryption, and \textit{fully} homormorphic encryption in particular, was envisioned for thirty years before the first lattice-based candidate was proposed in 2009 \cite{Gen09}.
    \item In practice, computational operations in lattice-based cryptography amount to matrix and vector operations modulo small integers, potentially over structured domains. Such operations are simple and highly parallelisable, allowing for relatively high performance compared to other post-quantum candidates. 
    \item Courtesy of Ajtai \cite{ajt96}, there exist reductions from average-case cryptographically relevant problems to worst-case instances over some related lattice problems. Such reductions allow cryptographic constructions that can be shown to be infeasible to break unless all instances over the corresponding lattice problem are easy to solve. Such properties are not commonly found for other classes of problems.
\end{itemize}

The central object of study is a lattice, which is an infinite, discrete set of regularly spaced points over $\RR^\nd$. These points are called \textit{lattice points}. 
\begin{definition}[Lattice]
    A lattice $\LL$ is a discrete additive subgroup of $\RR^\nd$.
\end{definition} 
The computational problems we discuss are usually concerned with finding something small in this lattice, such as the lattice point closest to the origin, or the lattice point closest to a given arbitrary point in $\RR^\nd$. In practice, the set of lattice points is parameterised by a \textit{lattice basis}. The choice of basis influences the hardness of solving the given computational problem. The cryptographic assumptions we are concerned with involve the creation of a lattice, which means solving the cryptographic problem also implies a solution (reduces) to an underlying lattice problem. We will see examples of this over different kinds of lattices throughout the text.

\subsection{Organisation}
This text is organised as follows. Section~\ref{sec:background} covers some theory of lattices, and proceeds to define four computational problems which lattice-based assumptions we discuss reduce to. Section~\ref{sec:LP} covers the assumptions themselves, namely Learning With Errors (LWE) and variants and Short Integers Solutions (SIS). The connections between these assumptions is also discussed. Section~\ref{sec:enc} focuses on encryption schemes from LWE, gradually giving more advanced schemes, such as homormorphic encryption and the now standardised ML-KEM (Kyber). As another major class of primitives, Section~\ref{sec:sig} presents a few digital signature constructions from lattices. Finally, Section~\ref{sec:further} is an epilogue discussing alternative assumptions in lattice cryptography.